% Acta de votación telemática - Consejo CEIC
% Movilización y paralización convocada por la FECh para el 18 de noviembre de 2026.
%
% Compilación:
%   Versión final:     pdflatex Acta_CEIC_Votacion_Movilizacion_18_de_Noviembre.tex   (dos veces)
%   Versión revisión:  pdflatex -jobname=Acta_CEIC_Votacion_Movilizacion_18_de_Noviembre_REVISION "\def\revision{}% Acta de votación telemática - Consejo CEIC
% Movilización y paralización convocada por la FECh para el 18 de noviembre de 2026.
%
% Compilación:
%   Versión final:     pdflatex Acta_CEIC_Votacion_Movilizacion_18_de_Noviembre.tex   (dos veces)
%   Versión revisión:  pdflatex -jobname=Acta_CEIC_Votacion_Movilizacion_18_de_Noviembre_REVISION "\def\revision{}% Acta de votación telemática - Consejo CEIC
% Movilización y paralización convocada por la FECh para el 18 de noviembre de 2026.
%
% Compilación:
%   Versión final:     pdflatex Acta_CEIC_Votacion_Movilizacion_18_de_Noviembre.tex   (dos veces)
%   Versión revisión:  pdflatex -jobname=Acta_CEIC_Votacion_Movilizacion_18_de_Noviembre_REVISION "\def\revision{}% Acta de votación telemática - Consejo CEIC
% Movilización y paralización convocada por la FECh para el 18 de noviembre de 2026.
%
% Compilación:
%   Versión final:     pdflatex Acta_CEIC_Votacion_Movilizacion_18_de_Noviembre.tex   (dos veces)
%   Versión revisión:  pdflatex -jobname=Acta_CEIC_Votacion_Movilizacion_18_de_Noviembre_REVISION "\def\revision{}\input{Acta_CEIC_Votacion_Movilizacion_18_de_Noviembre.tex}"
%
% En la versión de revisión, todo lo marcado con \dato{...} aparece en color:
% son datos que deben verificarse o completarse antes de publicar el acta.

\documentclass[12pt,letterpaper]{article}

\usepackage[T1]{fontenc}
\usepackage[utf8]{inputenc}
\usepackage[spanish,es-nolists,es-nodecimaldot]{babel}
\usepackage[letterpaper,left=1in,right=1in,top=1.35in,bottom=1.1in,headheight=48pt,headsep=24pt,footskip=36pt]{geometry}
\usepackage{graphicx}
\usepackage{amssymb}
\usepackage{xcolor}
\usepackage{fancyhdr}
\usepackage{lastpage}
\usepackage{titlesec}
\usepackage{enumitem}
\usepackage{textcomp}
\usepackage{array}
\usepackage{booktabs}

% ---------------------------------------------------------------------------
% Datos a verificar
% ---------------------------------------------------------------------------
\newif\ifrevision
\ifdefined\revision \revisiontrue \else \revisionfalse \fi
\newcommand{\dato}[1]{\ifrevision\textcolor[HTML]{C0392B}{#1}\else#1\fi}

% ---------------------------------------------------------------------------
% Datos del acta
% ---------------------------------------------------------------------------
\newcommand{\numeroacta}{\dato{[N°]}}
\newcommand{\fechaacta}{\dato{[dd-mm-2026]}}

% ---------------------------------------------------------------------------
% Formato
% ---------------------------------------------------------------------------
\setlength{\parindent}{1.5em}
\setlength{\parskip}{6pt}

\pagestyle{fancy}
\fancyhf{}
\fancyhead[L]{\raisebox{-0.35\height}{\includegraphics[height=1.15cm]{logo_ceic.png}}}
\fancyhead[R]{Acta N\textdegree\numeroacta\\\fechaacta}
\fancyfoot[C]{Página \thepage\ de \pageref{LastPage}}
\renewcommand{\headrulewidth}{0.4pt}
\renewcommand{\footrulewidth}{0.4pt}

\titleformat{\section}{\normalfont\Large\bfseries}{\thesection.}{0.35em}{}
\titleformat{\subsection}{\normalfont\large\bfseries}{\thesubsection.}{0.35em}{}
\titlespacing*{\section}{0pt}{3.0ex plus 1ex minus .2ex}{1.6ex plus .2ex}
\titlespacing*{\subsection}{0pt}{2.4ex plus 1ex minus .2ex}{1.2ex plus .2ex}

\renewcommand{\labelitemi}{\tiny$\blacksquare$}
\setlist[itemize]{leftmargin=2.2em,itemsep=2pt}
\setlist[enumerate]{leftmargin=2.2em,itemsep=2pt}

\newcolumntype{L}[1]{>{\raggedright\arraybackslash}p{#1}}
\newcolumntype{C}[1]{>{\centering\arraybackslash}p{#1}}

\begin{document}

% ---------------------------------------------------------------------------
% Encabezado
% ---------------------------------------------------------------------------
\begin{center}
  {\Large\bfseries Centro de Estudiantes de Ingeniería Comercial}\\
  Universidad de Chile, Facultad de Economía y Negocios

  \vspace{1.2em}
  \textbf{Acta de Votación Telemática del Consejo}
\end{center}

\vspace{0.8em}

\textbf{Materia:} Movilización y paralización de actividades académicas convocada por la FECh para el 18 de noviembre de 2026.

\textbf{Modalidad:} Votación telemática mediante formulario de Google Forms.

\textbf{Apertura de la votación:} \dato{[fecha y hora]}

\textbf{Cierre de la votación:} \dato{[fecha y hora]}

\vspace{1em}

\textbf{Tabla:}

\begin{enumerate}[itemsep=6pt]
  \item Antecedentes
  \item Modalidad de votación
  \item Mociones sometidas a votación
  \item Resultados
  \item Resolución
\end{enumerate}

% ---------------------------------------------------------------------------
\section{Antecedentes}

En el Pleno de la Federación de Estudiantes de la Universidad de Chile (FECh), celebrado el \dato{[fecha del pleno]}, se aprobó una moción de movilización y paralización de actividades académicas para el día 18 de noviembre de 2026.

En cumplimiento del compromiso asumido por la Mesa Directiva de presentar ante el Consejo este tipo de mociones, se convocó a una votación telemática con el fin de conocer la postura del Consejo respecto de la convocatoria.

Se deja constancia de que el 18 de noviembre coincide con el inicio del período de exámenes de la Universidad, por lo que no corresponde a una jornada de clases regulares, sino al comienzo del período destinado a la rendición de evaluaciones. Sin perjuicio de ello, se consideró pertinente someter la moción a votación del Consejo.

% ---------------------------------------------------------------------------
\section{Modalidad de votación}

La votación se realizó de forma telemática, mediante un formulario de Google Forms distribuido a los integrantes del Consejo. En cada moción, los participantes debían seleccionar una única alternativa: \emph{A favor}, \emph{En contra} o \emph{Abstención}.

De acuerdo con los estatutos, esta votación corresponde a un voto completo del Consejo. En consecuencia, solo se contabilizaron los votos de las delegaturas, de la Mesa Directiva y de las consejerías de escuela. Las consejerías temáticas (Género, Diversidad Sexual y Medio Ambiente) no tienen ponderación en este tipo de votación.

El resultado de cada estamento se calcula sobre los votos emitidos por sus integrantes y se pondera de la siguiente forma:

\begin{itemize}
  \item \textbf{Delegatura Mechona:} 20\%
  \item \textbf{Delegatura Plan Común:} 20\%
  \item \textbf{Delegatura Mención Economía:} 20\%
  \item \textbf{Delegatura Mención Administración:} 20\%
  \item \textbf{Mesa Directiva:} 10\%
  \item \textbf{Consejerías de Escuela:} 10\%
\end{itemize}

% ---------------------------------------------------------------------------
\section{Mociones sometidas a votación}

Se sometieron a votación dos mociones independientes, cuyos resultados se registran por separado:

\begin{enumerate}
  \item \textbf{Moción 1:} Aprobar la adhesión a la movilización convocada por la FECh para el día 18 de noviembre.

  Esta moción tuvo por objeto determinar si el Consejo respaldaba la adhesión a la movilización convocada por la Federación.

  \item \textbf{Moción 2:} Aprobar la paralización de actividades académicas para el día 18 de noviembre, en el marco de la movilización convocada por la FECh.

  Esta moción tuvo por objeto determinar si el Consejo respaldaba específicamente la paralización de actividades académicas durante dicha jornada.
\end{enumerate}

% ---------------------------------------------------------------------------
\section{Resultados}

\subsection{Registro de votos}

Se registraron los siguientes votos. Los integrantes que no figuran en la tabla no registraron voto en el formulario.

\begin{center}
\small\setlength{\tabcolsep}{4pt}
\begin{tabular}{L{4.3cm} L{6.2cm} C{2.2cm} C{2.2cm}}
  \toprule
  \textbf{Nombre} & \textbf{Estamento} & \textbf{Moción 1} & \textbf{Moción 2} \\
  \midrule
  Anabelee Araneda      & Delegatura Mechona                & A favor   & A favor   \\
  Antonio Muñoz         & Delegatura Mechona                & A favor   & A favor   \\
  Arturo Mora           & Delegatura Mechona                & En contra & En contra \\
  Juan Budge            & Delegatura Mechona                & En contra & En contra \\
  \midrule
  Martín Farías         & Delegatura Plan Común             & En contra & En contra \\
  \midrule
  Diego Peñailillo      & Delegatura Mención Economía       & A favor   & A favor   \\
  Nicolás Sanzana       & Delegatura Mención Economía       & A favor   & A favor   \\
  \midrule
  Sofía Cisternas       & Delegatura Mención Administración & En contra & En contra \\
  Juan José Esbeile     & Delegatura Mención Administración & En contra & En contra \\
  Diego Hernández       & Delegatura Mención Administración & A favor   & En contra \\
  Benjamín Gonzáles     & Delegatura Mención Administración & En contra & En contra \\
  \midrule
  Alan Mc Coll          & Mesa Directiva                    & En contra & En contra \\
  \midrule
  Benjamín Caro         & Consejerías de Escuela            & En contra & En contra \\
  \bottomrule
\end{tabular}
\end{center}

No registraron voto: Florencia Montiglio y José Toledo (Delegatura Plan Común); Juan José Tank y Nicolás Mella (Delegatura Mención Economía); Vicente Burgos (Consejerías de Escuela).

\subsection{Moción 1: Adhesión a la movilización convocada por la FECh}

\begin{center}
\small\setlength{\tabcolsep}{4pt}
\begin{tabular}{L{6.3cm} C{2.4cm} C{1.8cm} C{1.9cm} C{2.2cm}}
  \toprule
  \textbf{Estamento} & \textbf{Ponderación} & \textbf{A favor} & \textbf{En contra} & \textbf{Abstención} \\
  \midrule
  Delegatura Mechona                & 20\% & 50\%  & 50\%  & 0\% \\
  Delegatura Plan Común             & 20\% & 0\%   & 100\% & 0\% \\
  Delegatura Mención Economía       & 20\% & 100\% & 0\%   & 0\% \\
  Delegatura Mención Administración & 20\% & 25\%  & 75\%  & 0\% \\
  Mesa Directiva                    & 10\% & 0\%   & 100\% & 0\% \\
  Consejerías de Escuela            & 10\% & 0\%   & 100\% & 0\% \\
  \midrule
  \textbf{Resultado ponderado}      & \textbf{100\%} & \textbf{35\%} & \textbf{65\%} & \textbf{0\%} \\
  \bottomrule
\end{tabular}
\end{center}

\subsection{Moción 2: Paralización de actividades académicas}

\begin{center}
\small\setlength{\tabcolsep}{4pt}
\begin{tabular}{L{6.3cm} C{2.4cm} C{1.8cm} C{1.9cm} C{2.2cm}}
  \toprule
  \textbf{Estamento} & \textbf{Ponderación} & \textbf{A favor} & \textbf{En contra} & \textbf{Abstención} \\
  \midrule
  Delegatura Mechona                & 20\% & 50\%  & 50\%  & 0\% \\
  Delegatura Plan Común             & 20\% & 0\%   & 100\% & 0\% \\
  Delegatura Mención Economía       & 20\% & 100\% & 0\%   & 0\% \\
  Delegatura Mención Administración & 20\% & 0\%   & 100\% & 0\% \\
  Mesa Directiva                    & 10\% & 0\%   & 100\% & 0\% \\
  Consejerías de Escuela            & 10\% & 0\%   & 100\% & 0\% \\
  \midrule
  \textbf{Resultado ponderado}      & \textbf{100\%} & \textbf{30\%} & \textbf{70\%} & \textbf{0\%} \\
  \bottomrule
\end{tabular}
\end{center}

% ---------------------------------------------------------------------------
\section{Resolución}

En virtud de los resultados expuestos, el Consejo resuelve:

\begin{enumerate}
  \item \textbf{Moción 1:} Se \textbf{rechaza} la adhesión a la movilización convocada por la FECh para el día 18 de noviembre, con un 35\% de los votos ponderados a favor y un 65\% en contra.
  \item \textbf{Moción 2:} Se \textbf{rechaza} la paralización de actividades académicas para el día 18 de noviembre, con un 30\% de los votos ponderados a favor y un 70\% en contra.
\end{enumerate}

\vspace{0.5em}

Habiéndose computado la totalidad de los votos recibidos, se da por cerrada la votación.

\vspace{2.5em}

\noindent\textbf{Firma:}

\vspace{3em}

\noindent Alan Evan Mc Coll Albuja\\
Presidente CEIC

\end{document}
"
%
% En la versión de revisión, todo lo marcado con \dato{...} aparece en color:
% son datos que deben verificarse o completarse antes de publicar el acta.

\documentclass[12pt,letterpaper]{article}

\usepackage[T1]{fontenc}
\usepackage[utf8]{inputenc}
\usepackage[spanish,es-nolists,es-nodecimaldot]{babel}
\usepackage[letterpaper,left=1in,right=1in,top=1.35in,bottom=1.1in,headheight=48pt,headsep=24pt,footskip=36pt]{geometry}
\usepackage{graphicx}
\usepackage{amssymb}
\usepackage{xcolor}
\usepackage{fancyhdr}
\usepackage{lastpage}
\usepackage{titlesec}
\usepackage{enumitem}
\usepackage{textcomp}
\usepackage{array}
\usepackage{booktabs}

% ---------------------------------------------------------------------------
% Datos a verificar
% ---------------------------------------------------------------------------
\newif\ifrevision
\ifdefined\revision \revisiontrue \else \revisionfalse \fi
\newcommand{\dato}[1]{\ifrevision\textcolor[HTML]{C0392B}{#1}\else#1\fi}

% ---------------------------------------------------------------------------
% Datos del acta
% ---------------------------------------------------------------------------
\newcommand{\numeroacta}{\dato{[N°]}}
\newcommand{\fechaacta}{\dato{[dd-mm-2026]}}

% ---------------------------------------------------------------------------
% Formato
% ---------------------------------------------------------------------------
\setlength{\parindent}{1.5em}
\setlength{\parskip}{6pt}

\pagestyle{fancy}
\fancyhf{}
\fancyhead[L]{\raisebox{-0.35\height}{\includegraphics[height=1.15cm]{logo_ceic.png}}}
\fancyhead[R]{Acta N\textdegree\numeroacta\\\fechaacta}
\fancyfoot[C]{Página \thepage\ de \pageref{LastPage}}
\renewcommand{\headrulewidth}{0.4pt}
\renewcommand{\footrulewidth}{0.4pt}

\titleformat{\section}{\normalfont\Large\bfseries}{\thesection.}{0.35em}{}
\titleformat{\subsection}{\normalfont\large\bfseries}{\thesubsection.}{0.35em}{}
\titlespacing*{\section}{0pt}{3.0ex plus 1ex minus .2ex}{1.6ex plus .2ex}
\titlespacing*{\subsection}{0pt}{2.4ex plus 1ex minus .2ex}{1.2ex plus .2ex}

\renewcommand{\labelitemi}{\tiny$\blacksquare$}
\setlist[itemize]{leftmargin=2.2em,itemsep=2pt}
\setlist[enumerate]{leftmargin=2.2em,itemsep=2pt}

\newcolumntype{L}[1]{>{\raggedright\arraybackslash}p{#1}}
\newcolumntype{C}[1]{>{\centering\arraybackslash}p{#1}}

\begin{document}

% ---------------------------------------------------------------------------
% Encabezado
% ---------------------------------------------------------------------------
\begin{center}
  {\Large\bfseries Centro de Estudiantes de Ingeniería Comercial}\\
  Universidad de Chile, Facultad de Economía y Negocios

  \vspace{1.2em}
  \textbf{Acta de Votación Telemática del Consejo}
\end{center}

\vspace{0.8em}

\textbf{Materia:} Movilización y paralización de actividades académicas convocada por la FECh para el 18 de noviembre de 2026.

\textbf{Modalidad:} Votación telemática mediante formulario de Google Forms.

\textbf{Apertura de la votación:} \dato{[fecha y hora]}

\textbf{Cierre de la votación:} \dato{[fecha y hora]}

\vspace{1em}

\textbf{Tabla:}

\begin{enumerate}[itemsep=6pt]
  \item Antecedentes
  \item Modalidad de votación
  \item Mociones sometidas a votación
  \item Resultados
  \item Resolución
\end{enumerate}

% ---------------------------------------------------------------------------
\section{Antecedentes}

En el Pleno de la Federación de Estudiantes de la Universidad de Chile (FECh), celebrado el \dato{[fecha del pleno]}, se aprobó una moción de movilización y paralización de actividades académicas para el día 18 de noviembre de 2026.

En cumplimiento del compromiso asumido por la Mesa Directiva de presentar ante el Consejo este tipo de mociones, se convocó a una votación telemática con el fin de conocer la postura del Consejo respecto de la convocatoria.

Se deja constancia de que el 18 de noviembre coincide con el inicio del período de exámenes de la Universidad, por lo que no corresponde a una jornada de clases regulares, sino al comienzo del período destinado a la rendición de evaluaciones. Sin perjuicio de ello, se consideró pertinente someter la moción a votación del Consejo.

% ---------------------------------------------------------------------------
\section{Modalidad de votación}

La votación se realizó de forma telemática, mediante un formulario de Google Forms distribuido a los integrantes del Consejo. En cada moción, los participantes debían seleccionar una única alternativa: \emph{A favor}, \emph{En contra} o \emph{Abstención}.

De acuerdo con los estatutos, esta votación corresponde a un voto completo del Consejo. En consecuencia, solo se contabilizaron los votos de las delegaturas, de la Mesa Directiva y de las consejerías de escuela. Las consejerías temáticas (Género, Diversidad Sexual y Medio Ambiente) no tienen ponderación en este tipo de votación.

El resultado de cada estamento se calcula sobre los votos emitidos por sus integrantes y se pondera de la siguiente forma:

\begin{itemize}
  \item \textbf{Delegatura Mechona:} 20\%
  \item \textbf{Delegatura Plan Común:} 20\%
  \item \textbf{Delegatura Mención Economía:} 20\%
  \item \textbf{Delegatura Mención Administración:} 20\%
  \item \textbf{Mesa Directiva:} 10\%
  \item \textbf{Consejerías de Escuela:} 10\%
\end{itemize}

% ---------------------------------------------------------------------------
\section{Mociones sometidas a votación}

Se sometieron a votación dos mociones independientes, cuyos resultados se registran por separado:

\begin{enumerate}
  \item \textbf{Moción 1:} Aprobar la adhesión a la movilización convocada por la FECh para el día 18 de noviembre.

  Esta moción tuvo por objeto determinar si el Consejo respaldaba la adhesión a la movilización convocada por la Federación.

  \item \textbf{Moción 2:} Aprobar la paralización de actividades académicas para el día 18 de noviembre, en el marco de la movilización convocada por la FECh.

  Esta moción tuvo por objeto determinar si el Consejo respaldaba específicamente la paralización de actividades académicas durante dicha jornada.
\end{enumerate}

% ---------------------------------------------------------------------------
\section{Resultados}

\subsection{Registro de votos}

Se registraron los siguientes votos. Los integrantes que no figuran en la tabla no registraron voto en el formulario.

\begin{center}
\small\setlength{\tabcolsep}{4pt}
\begin{tabular}{L{4.3cm} L{6.2cm} C{2.2cm} C{2.2cm}}
  \toprule
  \textbf{Nombre} & \textbf{Estamento} & \textbf{Moción 1} & \textbf{Moción 2} \\
  \midrule
  Anabelee Araneda      & Delegatura Mechona                & A favor   & A favor   \\
  Antonio Muñoz         & Delegatura Mechona                & A favor   & A favor   \\
  Arturo Mora           & Delegatura Mechona                & En contra & En contra \\
  Juan Budge            & Delegatura Mechona                & En contra & En contra \\
  \midrule
  Martín Farías         & Delegatura Plan Común             & En contra & En contra \\
  \midrule
  Diego Peñailillo      & Delegatura Mención Economía       & A favor   & A favor   \\
  Nicolás Sanzana       & Delegatura Mención Economía       & A favor   & A favor   \\
  \midrule
  Sofía Cisternas       & Delegatura Mención Administración & En contra & En contra \\
  Juan José Esbeile     & Delegatura Mención Administración & En contra & En contra \\
  Diego Hernández       & Delegatura Mención Administración & A favor   & En contra \\
  Benjamín Gonzáles     & Delegatura Mención Administración & En contra & En contra \\
  \midrule
  Alan Mc Coll          & Mesa Directiva                    & En contra & En contra \\
  \midrule
  Benjamín Caro         & Consejerías de Escuela            & En contra & En contra \\
  \bottomrule
\end{tabular}
\end{center}

No registraron voto: Florencia Montiglio y José Toledo (Delegatura Plan Común); Juan José Tank y Nicolás Mella (Delegatura Mención Economía); Vicente Burgos (Consejerías de Escuela).

\subsection{Moción 1: Adhesión a la movilización convocada por la FECh}

\begin{center}
\small\setlength{\tabcolsep}{4pt}
\begin{tabular}{L{6.3cm} C{2.4cm} C{1.8cm} C{1.9cm} C{2.2cm}}
  \toprule
  \textbf{Estamento} & \textbf{Ponderación} & \textbf{A favor} & \textbf{En contra} & \textbf{Abstención} \\
  \midrule
  Delegatura Mechona                & 20\% & 50\%  & 50\%  & 0\% \\
  Delegatura Plan Común             & 20\% & 0\%   & 100\% & 0\% \\
  Delegatura Mención Economía       & 20\% & 100\% & 0\%   & 0\% \\
  Delegatura Mención Administración & 20\% & 25\%  & 75\%  & 0\% \\
  Mesa Directiva                    & 10\% & 0\%   & 100\% & 0\% \\
  Consejerías de Escuela            & 10\% & 0\%   & 100\% & 0\% \\
  \midrule
  \textbf{Resultado ponderado}      & \textbf{100\%} & \textbf{35\%} & \textbf{65\%} & \textbf{0\%} \\
  \bottomrule
\end{tabular}
\end{center}

\subsection{Moción 2: Paralización de actividades académicas}

\begin{center}
\small\setlength{\tabcolsep}{4pt}
\begin{tabular}{L{6.3cm} C{2.4cm} C{1.8cm} C{1.9cm} C{2.2cm}}
  \toprule
  \textbf{Estamento} & \textbf{Ponderación} & \textbf{A favor} & \textbf{En contra} & \textbf{Abstención} \\
  \midrule
  Delegatura Mechona                & 20\% & 50\%  & 50\%  & 0\% \\
  Delegatura Plan Común             & 20\% & 0\%   & 100\% & 0\% \\
  Delegatura Mención Economía       & 20\% & 100\% & 0\%   & 0\% \\
  Delegatura Mención Administración & 20\% & 0\%   & 100\% & 0\% \\
  Mesa Directiva                    & 10\% & 0\%   & 100\% & 0\% \\
  Consejerías de Escuela            & 10\% & 0\%   & 100\% & 0\% \\
  \midrule
  \textbf{Resultado ponderado}      & \textbf{100\%} & \textbf{30\%} & \textbf{70\%} & \textbf{0\%} \\
  \bottomrule
\end{tabular}
\end{center}

% ---------------------------------------------------------------------------
\section{Resolución}

En virtud de los resultados expuestos, el Consejo resuelve:

\begin{enumerate}
  \item \textbf{Moción 1:} Se \textbf{rechaza} la adhesión a la movilización convocada por la FECh para el día 18 de noviembre, con un 35\% de los votos ponderados a favor y un 65\% en contra.
  \item \textbf{Moción 2:} Se \textbf{rechaza} la paralización de actividades académicas para el día 18 de noviembre, con un 30\% de los votos ponderados a favor y un 70\% en contra.
\end{enumerate}

\vspace{0.5em}

Habiéndose computado la totalidad de los votos recibidos, se da por cerrada la votación.

\vspace{2.5em}

\noindent\textbf{Firma:}

\vspace{3em}

\noindent Alan Evan Mc Coll Albuja\\
Presidente CEIC

\end{document}
"
%
% En la versión de revisión, todo lo marcado con \dato{...} aparece en color:
% son datos que deben verificarse o completarse antes de publicar el acta.

\documentclass[12pt,letterpaper]{article}

\usepackage[T1]{fontenc}
\usepackage[utf8]{inputenc}
\usepackage[spanish,es-nolists,es-nodecimaldot]{babel}
\usepackage[letterpaper,left=1in,right=1in,top=1.35in,bottom=1.1in,headheight=48pt,headsep=24pt,footskip=36pt]{geometry}
\usepackage{graphicx}
\usepackage{amssymb}
\usepackage{xcolor}
\usepackage{fancyhdr}
\usepackage{lastpage}
\usepackage{titlesec}
\usepackage{enumitem}
\usepackage{textcomp}
\usepackage{array}
\usepackage{booktabs}

% ---------------------------------------------------------------------------
% Datos a verificar
% ---------------------------------------------------------------------------
\newif\ifrevision
\ifdefined\revision \revisiontrue \else \revisionfalse \fi
\newcommand{\dato}[1]{\ifrevision\textcolor[HTML]{C0392B}{#1}\else#1\fi}

% ---------------------------------------------------------------------------
% Datos del acta
% ---------------------------------------------------------------------------
\newcommand{\numeroacta}{\dato{[N°]}}
\newcommand{\fechaacta}{\dato{[dd-mm-2026]}}

% ---------------------------------------------------------------------------
% Formato
% ---------------------------------------------------------------------------
\setlength{\parindent}{1.5em}
\setlength{\parskip}{6pt}

\pagestyle{fancy}
\fancyhf{}
\fancyhead[L]{\raisebox{-0.35\height}{\includegraphics[height=1.15cm]{logo_ceic.png}}}
\fancyhead[R]{Acta N\textdegree\numeroacta\\\fechaacta}
\fancyfoot[C]{Página \thepage\ de \pageref{LastPage}}
\renewcommand{\headrulewidth}{0.4pt}
\renewcommand{\footrulewidth}{0.4pt}

\titleformat{\section}{\normalfont\Large\bfseries}{\thesection.}{0.35em}{}
\titleformat{\subsection}{\normalfont\large\bfseries}{\thesubsection.}{0.35em}{}
\titlespacing*{\section}{0pt}{3.0ex plus 1ex minus .2ex}{1.6ex plus .2ex}
\titlespacing*{\subsection}{0pt}{2.4ex plus 1ex minus .2ex}{1.2ex plus .2ex}

\renewcommand{\labelitemi}{\tiny$\blacksquare$}
\setlist[itemize]{leftmargin=2.2em,itemsep=2pt}
\setlist[enumerate]{leftmargin=2.2em,itemsep=2pt}

\newcolumntype{L}[1]{>{\raggedright\arraybackslash}p{#1}}
\newcolumntype{C}[1]{>{\centering\arraybackslash}p{#1}}

\begin{document}

% ---------------------------------------------------------------------------
% Encabezado
% ---------------------------------------------------------------------------
\begin{center}
  {\Large\bfseries Centro de Estudiantes de Ingeniería Comercial}\\
  Universidad de Chile, Facultad de Economía y Negocios

  \vspace{1.2em}
  \textbf{Acta de Votación Telemática del Consejo}
\end{center}

\vspace{0.8em}

\textbf{Materia:} Movilización y paralización de actividades académicas convocada por la FECh para el 18 de noviembre de 2026.

\textbf{Modalidad:} Votación telemática mediante formulario de Google Forms.

\textbf{Apertura de la votación:} \dato{[fecha y hora]}

\textbf{Cierre de la votación:} \dato{[fecha y hora]}

\vspace{1em}

\textbf{Tabla:}

\begin{enumerate}[itemsep=6pt]
  \item Antecedentes
  \item Modalidad de votación
  \item Mociones sometidas a votación
  \item Resultados
  \item Resolución
\end{enumerate}

% ---------------------------------------------------------------------------
\section{Antecedentes}

En el Pleno de la Federación de Estudiantes de la Universidad de Chile (FECh), celebrado el \dato{[fecha del pleno]}, se aprobó una moción de movilización y paralización de actividades académicas para el día 18 de noviembre de 2026.

En cumplimiento del compromiso asumido por la Mesa Directiva de presentar ante el Consejo este tipo de mociones, se convocó a una votación telemática con el fin de conocer la postura del Consejo respecto de la convocatoria.

Se deja constancia de que el 18 de noviembre coincide con el inicio del período de exámenes de la Universidad, por lo que no corresponde a una jornada de clases regulares, sino al comienzo del período destinado a la rendición de evaluaciones. Sin perjuicio de ello, se consideró pertinente someter la moción a votación del Consejo.

% ---------------------------------------------------------------------------
\section{Modalidad de votación}

La votación se realizó de forma telemática, mediante un formulario de Google Forms distribuido a los integrantes del Consejo. En cada moción, los participantes debían seleccionar una única alternativa: \emph{A favor}, \emph{En contra} o \emph{Abstención}.

De acuerdo con los estatutos, esta votación corresponde a un voto completo del Consejo. En consecuencia, solo se contabilizaron los votos de las delegaturas, de la Mesa Directiva y de las consejerías de escuela. Las consejerías temáticas (Género, Diversidad Sexual y Medio Ambiente) no tienen ponderación en este tipo de votación.

El resultado de cada estamento se calcula sobre los votos emitidos por sus integrantes y se pondera de la siguiente forma:

\begin{itemize}
  \item \textbf{Delegatura Mechona:} 20\%
  \item \textbf{Delegatura Plan Común:} 20\%
  \item \textbf{Delegatura Mención Economía:} 20\%
  \item \textbf{Delegatura Mención Administración:} 20\%
  \item \textbf{Mesa Directiva:} 10\%
  \item \textbf{Consejerías de Escuela:} 10\%
\end{itemize}

% ---------------------------------------------------------------------------
\section{Mociones sometidas a votación}

Se sometieron a votación dos mociones independientes, cuyos resultados se registran por separado:

\begin{enumerate}
  \item \textbf{Moción 1:} Aprobar la adhesión a la movilización convocada por la FECh para el día 18 de noviembre.

  Esta moción tuvo por objeto determinar si el Consejo respaldaba la adhesión a la movilización convocada por la Federación.

  \item \textbf{Moción 2:} Aprobar la paralización de actividades académicas para el día 18 de noviembre, en el marco de la movilización convocada por la FECh.

  Esta moción tuvo por objeto determinar si el Consejo respaldaba específicamente la paralización de actividades académicas durante dicha jornada.
\end{enumerate}

% ---------------------------------------------------------------------------
\section{Resultados}

\subsection{Registro de votos}

Se registraron los siguientes votos. Los integrantes que no figuran en la tabla no registraron voto en el formulario.

\begin{center}
\small\setlength{\tabcolsep}{4pt}
\begin{tabular}{L{4.3cm} L{6.2cm} C{2.2cm} C{2.2cm}}
  \toprule
  \textbf{Nombre} & \textbf{Estamento} & \textbf{Moción 1} & \textbf{Moción 2} \\
  \midrule
  Anabelee Araneda      & Delegatura Mechona                & A favor   & A favor   \\
  Antonio Muñoz         & Delegatura Mechona                & A favor   & A favor   \\
  Arturo Mora           & Delegatura Mechona                & En contra & En contra \\
  Juan Budge            & Delegatura Mechona                & En contra & En contra \\
  \midrule
  Martín Farías         & Delegatura Plan Común             & En contra & En contra \\
  \midrule
  Diego Peñailillo      & Delegatura Mención Economía       & A favor   & A favor   \\
  Nicolás Sanzana       & Delegatura Mención Economía       & A favor   & A favor   \\
  \midrule
  Sofía Cisternas       & Delegatura Mención Administración & En contra & En contra \\
  Juan José Esbeile     & Delegatura Mención Administración & En contra & En contra \\
  Diego Hernández       & Delegatura Mención Administración & A favor   & En contra \\
  Benjamín Gonzáles     & Delegatura Mención Administración & En contra & En contra \\
  \midrule
  Alan Mc Coll          & Mesa Directiva                    & En contra & En contra \\
  \midrule
  Benjamín Caro         & Consejerías de Escuela            & En contra & En contra \\
  \bottomrule
\end{tabular}
\end{center}

No registraron voto: Florencia Montiglio y José Toledo (Delegatura Plan Común); Juan José Tank y Nicolás Mella (Delegatura Mención Economía); Vicente Burgos (Consejerías de Escuela).

\subsection{Moción 1: Adhesión a la movilización convocada por la FECh}

\begin{center}
\small\setlength{\tabcolsep}{4pt}
\begin{tabular}{L{6.3cm} C{2.4cm} C{1.8cm} C{1.9cm} C{2.2cm}}
  \toprule
  \textbf{Estamento} & \textbf{Ponderación} & \textbf{A favor} & \textbf{En contra} & \textbf{Abstención} \\
  \midrule
  Delegatura Mechona                & 20\% & 50\%  & 50\%  & 0\% \\
  Delegatura Plan Común             & 20\% & 0\%   & 100\% & 0\% \\
  Delegatura Mención Economía       & 20\% & 100\% & 0\%   & 0\% \\
  Delegatura Mención Administración & 20\% & 25\%  & 75\%  & 0\% \\
  Mesa Directiva                    & 10\% & 0\%   & 100\% & 0\% \\
  Consejerías de Escuela            & 10\% & 0\%   & 100\% & 0\% \\
  \midrule
  \textbf{Resultado ponderado}      & \textbf{100\%} & \textbf{35\%} & \textbf{65\%} & \textbf{0\%} \\
  \bottomrule
\end{tabular}
\end{center}

\subsection{Moción 2: Paralización de actividades académicas}

\begin{center}
\small\setlength{\tabcolsep}{4pt}
\begin{tabular}{L{6.3cm} C{2.4cm} C{1.8cm} C{1.9cm} C{2.2cm}}
  \toprule
  \textbf{Estamento} & \textbf{Ponderación} & \textbf{A favor} & \textbf{En contra} & \textbf{Abstención} \\
  \midrule
  Delegatura Mechona                & 20\% & 50\%  & 50\%  & 0\% \\
  Delegatura Plan Común             & 20\% & 0\%   & 100\% & 0\% \\
  Delegatura Mención Economía       & 20\% & 100\% & 0\%   & 0\% \\
  Delegatura Mención Administración & 20\% & 0\%   & 100\% & 0\% \\
  Mesa Directiva                    & 10\% & 0\%   & 100\% & 0\% \\
  Consejerías de Escuela            & 10\% & 0\%   & 100\% & 0\% \\
  \midrule
  \textbf{Resultado ponderado}      & \textbf{100\%} & \textbf{30\%} & \textbf{70\%} & \textbf{0\%} \\
  \bottomrule
\end{tabular}
\end{center}

% ---------------------------------------------------------------------------
\section{Resolución}

En virtud de los resultados expuestos, el Consejo resuelve:

\begin{enumerate}
  \item \textbf{Moción 1:} Se \textbf{rechaza} la adhesión a la movilización convocada por la FECh para el día 18 de noviembre, con un 35\% de los votos ponderados a favor y un 65\% en contra.
  \item \textbf{Moción 2:} Se \textbf{rechaza} la paralización de actividades académicas para el día 18 de noviembre, con un 30\% de los votos ponderados a favor y un 70\% en contra.
\end{enumerate}

\vspace{0.5em}

Habiéndose computado la totalidad de los votos recibidos, se da por cerrada la votación.

\vspace{2.5em}

\noindent\textbf{Firma:}

\vspace{3em}

\noindent Alan Evan Mc Coll Albuja\\
Presidente CEIC

\end{document}
"
%
% En la versión de revisión, todo lo marcado con \dato{...} aparece en color:
% son datos que deben verificarse o completarse antes de publicar el acta.

\documentclass[12pt,letterpaper]{article}

\usepackage[T1]{fontenc}
\usepackage[utf8]{inputenc}
\usepackage[spanish,es-nolists,es-nodecimaldot]{babel}
\usepackage[letterpaper,left=1in,right=1in,top=1.35in,bottom=1.1in,headheight=48pt,headsep=24pt,footskip=36pt]{geometry}
\usepackage{graphicx}
\usepackage{amssymb}
\usepackage{xcolor}
\usepackage{fancyhdr}
\usepackage{lastpage}
\usepackage{titlesec}
\usepackage{enumitem}
\usepackage{textcomp}
\usepackage{array}
\usepackage{booktabs}

% ---------------------------------------------------------------------------
% Datos a verificar
% ---------------------------------------------------------------------------
\newif\ifrevision
\ifdefined\revision \revisiontrue \else \revisionfalse \fi
\newcommand{\dato}[1]{\ifrevision\textcolor[HTML]{C0392B}{#1}\else#1\fi}

% ---------------------------------------------------------------------------
% Datos del acta
% ---------------------------------------------------------------------------
\newcommand{\numeroacta}{\dato{[N°]}}
\newcommand{\fechaacta}{\dato{[dd-mm-2026]}}

% ---------------------------------------------------------------------------
% Formato
% ---------------------------------------------------------------------------
\setlength{\parindent}{1.5em}
\setlength{\parskip}{6pt}

\pagestyle{fancy}
\fancyhf{}
\fancyhead[L]{\raisebox{-0.35\height}{\includegraphics[height=1.15cm]{logo_ceic.png}}}
\fancyhead[R]{Acta N\textdegree\numeroacta\\\fechaacta}
\fancyfoot[C]{Página \thepage\ de \pageref{LastPage}}
\renewcommand{\headrulewidth}{0.4pt}
\renewcommand{\footrulewidth}{0.4pt}

\titleformat{\section}{\normalfont\Large\bfseries}{\thesection.}{0.35em}{}
\titleformat{\subsection}{\normalfont\large\bfseries}{\thesubsection.}{0.35em}{}
\titlespacing*{\section}{0pt}{3.0ex plus 1ex minus .2ex}{1.6ex plus .2ex}
\titlespacing*{\subsection}{0pt}{2.4ex plus 1ex minus .2ex}{1.2ex plus .2ex}

\renewcommand{\labelitemi}{\tiny$\blacksquare$}
\setlist[itemize]{leftmargin=2.2em,itemsep=2pt}
\setlist[enumerate]{leftmargin=2.2em,itemsep=2pt}

\newcolumntype{L}[1]{>{\raggedright\arraybackslash}p{#1}}
\newcolumntype{C}[1]{>{\centering\arraybackslash}p{#1}}

\begin{document}

% ---------------------------------------------------------------------------
% Encabezado
% ---------------------------------------------------------------------------
\begin{center}
  {\Large\bfseries Centro de Estudiantes de Ingeniería Comercial}\\
  Universidad de Chile, Facultad de Economía y Negocios

  \vspace{1.2em}
  \textbf{Acta de Votación Telemática del Consejo}
\end{center}

\vspace{0.8em}

\textbf{Materia:} Movilización y paralización de actividades académicas convocada por la FECh para el 18 de noviembre de 2026.

\textbf{Modalidad:} Votación telemática mediante formulario de Google Forms.

\textbf{Apertura de la votación:} \dato{[fecha y hora]}

\textbf{Cierre de la votación:} \dato{[fecha y hora]}

\vspace{1em}

\textbf{Tabla:}

\begin{enumerate}[itemsep=6pt]
  \item Antecedentes
  \item Modalidad de votación
  \item Mociones sometidas a votación
  \item Resultados
  \item Resolución
\end{enumerate}

% ---------------------------------------------------------------------------
\section{Antecedentes}

En el Pleno de la Federación de Estudiantes de la Universidad de Chile (FECh), celebrado el \dato{[fecha del pleno]}, se aprobó una moción de movilización y paralización de actividades académicas para el día 18 de noviembre de 2026.

En cumplimiento del compromiso asumido por la Mesa Directiva de presentar ante el Consejo este tipo de mociones, se convocó a una votación telemática con el fin de conocer la postura del Consejo respecto de la convocatoria.

Se deja constancia de que el 18 de noviembre coincide con el inicio del período de exámenes de la Universidad, por lo que no corresponde a una jornada de clases regulares, sino al comienzo del período destinado a la rendición de evaluaciones. Sin perjuicio de ello, se consideró pertinente someter la moción a votación del Consejo.

% ---------------------------------------------------------------------------
\section{Modalidad de votación}

La votación se realizó de forma telemática, mediante un formulario de Google Forms distribuido a los integrantes del Consejo. En cada moción, los participantes debían seleccionar una única alternativa: \emph{A favor}, \emph{En contra} o \emph{Abstención}.

De acuerdo con los estatutos, esta votación corresponde a un voto completo del Consejo. En consecuencia, solo se contabilizaron los votos de las delegaturas, de la Mesa Directiva y de las consejerías de escuela. Las consejerías temáticas (Género, Diversidad Sexual y Medio Ambiente) no tienen ponderación en este tipo de votación.

El resultado de cada estamento se calcula sobre los votos emitidos por sus integrantes y se pondera de la siguiente forma:

\begin{itemize}
  \item \textbf{Delegatura Mechona:} 20\%
  \item \textbf{Delegatura Plan Común:} 20\%
  \item \textbf{Delegatura Mención Economía:} 20\%
  \item \textbf{Delegatura Mención Administración:} 20\%
  \item \textbf{Mesa Directiva:} 10\%
  \item \textbf{Consejerías de Escuela:} 10\%
\end{itemize}

% ---------------------------------------------------------------------------
\section{Mociones sometidas a votación}

Se sometieron a votación dos mociones independientes, cuyos resultados se registran por separado:

\begin{enumerate}
  \item \textbf{Moción 1:} Aprobar la adhesión a la movilización convocada por la FECh para el día 18 de noviembre.

  Esta moción tuvo por objeto determinar si el Consejo respaldaba la adhesión a la movilización convocada por la Federación.

  \item \textbf{Moción 2:} Aprobar la paralización de actividades académicas para el día 18 de noviembre, en el marco de la movilización convocada por la FECh.

  Esta moción tuvo por objeto determinar si el Consejo respaldaba específicamente la paralización de actividades académicas durante dicha jornada.
\end{enumerate}

% ---------------------------------------------------------------------------
\section{Resultados}

\subsection{Registro de votos}

Se registraron los siguientes votos. Los integrantes que no figuran en la tabla no registraron voto en el formulario.

\begin{center}
\small\setlength{\tabcolsep}{4pt}
\begin{tabular}{L{4.3cm} L{6.2cm} C{2.2cm} C{2.2cm}}
  \toprule
  \textbf{Nombre} & \textbf{Estamento} & \textbf{Moción 1} & \textbf{Moción 2} \\
  \midrule
  Anabelee Araneda      & Delegatura Mechona                & A favor   & A favor   \\
  Antonio Muñoz         & Delegatura Mechona                & A favor   & A favor   \\
  Arturo Mora           & Delegatura Mechona                & En contra & En contra \\
  Juan Budge            & Delegatura Mechona                & En contra & En contra \\
  \midrule
  Martín Farías         & Delegatura Plan Común             & En contra & En contra \\
  \midrule
  Diego Peñailillo      & Delegatura Mención Economía       & A favor   & A favor   \\
  Nicolás Sanzana       & Delegatura Mención Economía       & A favor   & A favor   \\
  \midrule
  Sofía Cisternas       & Delegatura Mención Administración & En contra & En contra \\
  Juan José Esbeile     & Delegatura Mención Administración & En contra & En contra \\
  Diego Hernández       & Delegatura Mención Administración & A favor   & En contra \\
  Benjamín Gonzáles     & Delegatura Mención Administración & En contra & En contra \\
  \midrule
  Alan Mc Coll          & Mesa Directiva                    & En contra & En contra \\
  \midrule
  Benjamín Caro         & Consejerías de Escuela            & En contra & En contra \\
  \bottomrule
\end{tabular}
\end{center}

No registraron voto: Florencia Montiglio y José Toledo (Delegatura Plan Común); Juan José Tank y Nicolás Mella (Delegatura Mención Economía); Vicente Burgos (Consejerías de Escuela).

\subsection{Moción 1: Adhesión a la movilización convocada por la FECh}

\begin{center}
\small\setlength{\tabcolsep}{4pt}
\begin{tabular}{L{6.3cm} C{2.4cm} C{1.8cm} C{1.9cm} C{2.2cm}}
  \toprule
  \textbf{Estamento} & \textbf{Ponderación} & \textbf{A favor} & \textbf{En contra} & \textbf{Abstención} \\
  \midrule
  Delegatura Mechona                & 20\% & 50\%  & 50\%  & 0\% \\
  Delegatura Plan Común             & 20\% & 0\%   & 100\% & 0\% \\
  Delegatura Mención Economía       & 20\% & 100\% & 0\%   & 0\% \\
  Delegatura Mención Administración & 20\% & 25\%  & 75\%  & 0\% \\
  Mesa Directiva                    & 10\% & 0\%   & 100\% & 0\% \\
  Consejerías de Escuela            & 10\% & 0\%   & 100\% & 0\% \\
  \midrule
  \textbf{Resultado ponderado}      & \textbf{100\%} & \textbf{35\%} & \textbf{65\%} & \textbf{0\%} \\
  \bottomrule
\end{tabular}
\end{center}

\subsection{Moción 2: Paralización de actividades académicas}

\begin{center}
\small\setlength{\tabcolsep}{4pt}
\begin{tabular}{L{6.3cm} C{2.4cm} C{1.8cm} C{1.9cm} C{2.2cm}}
  \toprule
  \textbf{Estamento} & \textbf{Ponderación} & \textbf{A favor} & \textbf{En contra} & \textbf{Abstención} \\
  \midrule
  Delegatura Mechona                & 20\% & 50\%  & 50\%  & 0\% \\
  Delegatura Plan Común             & 20\% & 0\%   & 100\% & 0\% \\
  Delegatura Mención Economía       & 20\% & 100\% & 0\%   & 0\% \\
  Delegatura Mención Administración & 20\% & 0\%   & 100\% & 0\% \\
  Mesa Directiva                    & 10\% & 0\%   & 100\% & 0\% \\
  Consejerías de Escuela            & 10\% & 0\%   & 100\% & 0\% \\
  \midrule
  \textbf{Resultado ponderado}      & \textbf{100\%} & \textbf{30\%} & \textbf{70\%} & \textbf{0\%} \\
  \bottomrule
\end{tabular}
\end{center}

% ---------------------------------------------------------------------------
\section{Resolución}

En virtud de los resultados expuestos, el Consejo resuelve:

\begin{enumerate}
  \item \textbf{Moción 1:} Se \textbf{rechaza} la adhesión a la movilización convocada por la FECh para el día 18 de noviembre, con un 35\% de los votos ponderados a favor y un 65\% en contra.
  \item \textbf{Moción 2:} Se \textbf{rechaza} la paralización de actividades académicas para el día 18 de noviembre, con un 30\% de los votos ponderados a favor y un 70\% en contra.
\end{enumerate}

\vspace{0.5em}

Habiéndose computado la totalidad de los votos recibidos, se da por cerrada la votación.

\vspace{2.5em}

\noindent\textbf{Firma:}

\vspace{3em}

\noindent Alan Evan Mc Coll Albuja\\
Presidente CEIC

\end{document}
